\documentclass[11pt]{amsart}

\usepackage[T1]{fontenc}
\usepackage{lmodern}
\usepackage{amssymb}
\usepackage{microtype}
\usepackage{booktabs}
\usepackage{xcolor}
\usepackage{url}
\usepackage{hyperref}

\makeatletter
\g@addto@macro\UrlBreaks{\do\_\do\.}
\makeatother

\setlength{\emergencystretch}{3em}
\hypersetup{
  colorlinks=true,
  linkcolor=blue!55!black,
  citecolor=green!40!black,
  urlcolor=blue!55!black,
  pdftitle={Lagrange Collisions and Cover Relations for Rational Dyck Paths},
  pdfauthor={Alex Chengyu Li}
}

\newtheorem{theorem}{Theorem}[section]
\newtheorem{proposition}[theorem]{Proposition}
\newtheorem{lemma}[theorem]{Lemma}
\newtheorem{corollary}[theorem]{Corollary}
\theoremstyle{definition}
\newtheorem{definition}[theorem]{Definition}
\newtheorem{remark}[theorem]{Remark}

\newcommand{\RR}{\mathrm{R}}
\newcommand{\UU}{\mathrm{U}}
\newcommand{\cD}{\mathcal D}
\newcommand{\Mat}{\operatorname{Mat}}
\newcommand{\tr}{\operatorname{tr}}
\newcommand{\Fib}{\operatorname{Fib}}
\newcommand{\Ftwo}{\mathbb F_2}

\title[Lagrange collisions and cover relations]
{Lagrange Collisions and Cover Relations for Rational Dyck Paths}

\author{Alex Chengyu Li}
\address{Independent Researcher, Tokyo, Japan; ORCID 0009-0008-4516-8946}
\email{cl7076@nyu.edu}
\date{August 2026}

\subjclass[2020]{Primary 05A15, 06A07; Secondary 11A55, 05C70}
\keywords{lattice paths, continuants, Lagrange spectrum, band graphs,
cover relations, algebraic combinatorics, exact certificates}

\begin{document}

\begin{abstract}
Schiffler associated matching and Lagrange orders to rational Dyck paths and
posed two problems about their equality and cover relations.  At every fixed
coprime endpoint, we classify the band-graph isomorphism classes: they are the
orbits of an explicit reversal involution and therefore have size at most two.
This gives a closed count of the graph fibres and reduces Lagrange-value
rigidity to collision between distinct fibres.  Such a collision occurs in
$\cD(17,9)$; two exact paths have the same Lagrange value but nonisomorphic
band graphs, and independent enumerations show that total length $26$ is
minimal.

The matching order also has intrinsically nonlocal cover geometry.  Its scores
at a fixed endpoint have one parity, its first $n$ levels in $\cD(n,n-1)$ are
explicit singletons, and a family of covers has Hamming distance $2n-6$.
Thus cover pairs cannot all differ in a uniformly bounded number of step
positions.  To recover the complete cover relation in both the matching and
Lagrange orders, we derive exact matrix intervals for rational-Dyck prefix
cylinders.  A best-first traversal emits complete score levels, while finite
prefix-antichain certificates verify that no intermediate level is omitted.
An exact four-case continuant formula separately describes every local
interchange.  All finite certificates use integer arithmetic or reduced
rational squares.
\end{abstract}

\maketitle

\section{The problems and main results}

Let $a>b>0$ be coprime.  The set $\cD(a,b)$ consists of the words with $a$
letters $\RR$ and $b$ letters $\UU$ for which every prefix $v$ satisfies
\[
  a\,\#_{\UU}(v)\leq b\,\#_{\RR}(v).
\]
Schiffler~\cite{Schiffler} associates two strict partial orders to these
paths.  The matching score is the numerator of a finite continued fraction,
and the Lagrange score is the Lagrange number of an associated periodic
continued fraction.  The underlying snake and band graphs belong to the graph
calculus developed by \c{C}anak\c{c}\i{} and Schiffler
\cite{CanakciSchifflerBand}.  Problem~6.3 of \cite{Schiffler} asks whether
equality of Lagrange values for two paths in the same $\cD(a,b)$ implies
isomorphism of their band graphs.  Problem 6.5 asks for a path description of
the cover relations in either order.  Apruzzese and Cong
\cite{ApruzzeseCong} prove the common unique maximum and the supremum of all
Lagrange values; they reiterate that comparison criteria alone do not exclude
an intermediate score level.

Two structural questions guide the paper.  First, how much of a rational-Dyck
path is remembered by its band graph and by the coarser Lagrange value?
Second, can adjacency in either score order be generated by bounded local
changes of the path?  We find rigidity up to an explicit involution at the
graph-fibre level but failure after passage to the scalar Lagrange value.  For
the matching order, explicit initial levels and a family of increasingly
distant covers show that cover pairs need not have bounded Hamming distance.
Exact prefix intervals then supply
a complete global certificate system for both orders.

The first structural point is that band-graph isomorphism itself has a simple
path description.  Reverse the cyclic step word of $\omega$ and rotate it to
the unique member of $\cD(a,b)$ in that rotation class; call the result
$\tau(\omega)$.  The existence and uniqueness are the coprime cycle lemma,
recalled in Section~\ref{sec:band-classes}.

\begin{theorem}[Band-graph classes]\label{thm:band-classes}
For $\omega,\omega'\in\cD(a,b)$,
\[
 \overline G(\omega)\cong\overline G(\omega')
 \quad\Longleftrightarrow\quad
 \omega'=\omega\ \hbox{or}\ \omega'=\tau(\omega).
\]
In particular, every band-graph isomorphism class has size at most two.
\end{theorem}

This makes the equality question concrete: one must decide whether a Lagrange
level can meet two different $\tau$-orbits.  Our first counterexample does so.

\begin{theorem}[Counterexample to Problem 6.3]\label{thm:counterexample}
There are paths $\omega,\omega'\in\cD(17,9)$ such that
\[
 L(\omega)=L(\omega')
 \quad\hbox{but}\quad
 \overline G(\omega)\not\cong\overline G(\omega').
\]
They can be chosen with step words
\begin{align*}
 \omega  &={\tt RRRRRRURRURRUURURRURRRUURU},\\
 \omega' &={\tt RRRRRRURUURURRURRRURRURRUU}.
\end{align*}
\end{theorem}

The equality and nonisomorphism are finite exact statements and do not depend
on floating-point approximations.

\paragraph{Concurrent work.}
The first public version of this manuscript was deposited on 17 August 2026
under version DOI \url{https://doi.org/10.5281/zenodo.21980890}.  Tang and
Zhang posted independent work on 19 August 2026 giving a different pair in
$\cD(17,9)$ with a different common Lagrange value~\cite{TangZhang}.  We cite
their paper as a concurrent negative answer to Problem~6.3.  The band-graph
classification, the length-$25$ minimality result, and the results on the
cover relations in Problem~6.5 are separate contributions of the present
paper.

Our answer to the cover problem is constructive and global.  Explicit lower
and upper score bounds on every path-prefix cylinder yield both a score-ordered
generator and a locally checkable prefix-antichain certificate for any
proposed cover.  No bounded catalogue of local swaps is required.

\begin{theorem}[Constructive cover characterization]\label{thm:global-cover}
For every coprime $a>b>0$, there is an exact best-first traversal of the
rational-Dyck prefix tree that emits the complete matching-score levels, and
also the complete Lagrange-score levels, in strictly increasing batches.
For either order, the cover relation is the union of the complete bipartite
relations between consecutive emitted batches.

Equivalently, paths $x<y$ form a cover if and only if there is a finite prefix
antichain partitioning $\cD(a,b)$ such that the score interval of every
antichain cylinder is certified to lie at or below the score of $x$, or at or
above the score of $y$.
\end{theorem}

The bounds and the certificate criterion are given explicitly in
Section~\ref{sec:covers}.  The matching order also admits useful uniform
structure.

\begin{theorem}[Endpoint parity]\label{thm:parity}
For every $\omega\in\cD(a,b)$,
\[
 M(\omega)\equiv0\pmod2
 \quad\Longleftrightarrow\quad
 3\mid(a-b).
\]
Consequently all matching scores at a fixed endpoint have the same parity,
and any two realized scores differing by two are consecutive levels.
\end{theorem}

Finally, the next theorem gives a closed initial segment and shows why the
global form of Theorem~\ref{thm:global-cover} is necessary.  For a path
$\omega$, write $A_E(\omega)$ for its word in the adjacency blocks $D,E$ of
Section~\ref{sec:matrices}.

\begin{theorem}[Initial matching levels and unbounded nonlocality]
\label{thm:nonlocal}
For every $n\geq4$, the paths
\[
 X_n=\RR\RR\RR(\UU\RR)^{n-3}\UU\UU,
 \qquad
 Y_n=\RR\RR(\UU\RR)^{n-4}\RR\UU\UU\RR\UU
\]
form a matching cover in $\cD(n,n-1)$ and have Hamming distance $2n-6$.
More precisely, the first $n$ distinct matching levels are singletons.  The
minimum is the path with block word $ED^{2n-3}$; the next $n-2$ paths have
block words
\[
 ED^{2r}ED^{2n-5-2r}E,
 \qquad 0\leq r\leq n-3,
\]
listed in decreasing order of $r$; and $Y_n$ is the unique path at the next
level.  In this list $X_n$ is the level immediately below $Y_n$.
\end{theorem}

Theorem~\ref{thm:band-classes} and the parity and initial-level theorems are
uniform structural results.  Their proofs use, respectively, the coprime cycle
lemma, a two-by-two calculation over $\Ftwo$, and an elementary
Pell-continuant extremal lemma.

\section{Background and exact matrix model}\label{sec:matrices}

For a path $\omega$, form a finite coefficient word $A(\omega)$ from adjacent
steps: an equal pair contributes the two coefficients $1,1$, while an unequal
pair contributes the coefficient $2$.  Schiffler's matching score $M(\omega)$
is the numerator of the continued fraction with coefficient word
$A(\omega)$.  Closing the step word cyclically adds one unequal edge, and the
corresponding periodic continued fraction has Lagrange number $L(\omega)$.
The two strict orders considered here are the pullbacks of the usual order on
these scores:
\[
 \omega<_M\eta\iff M(\omega)<M(\eta),\qquad
 \omega<_L\eta\iff L(\omega)<L(\eta).
\]
Paths with equal scores are therefore incomparable.  A cover is a pair of
paths in two consecutive distinct realized score levels.

Put
\[
 T(c)=\begin{pmatrix}c&1\\1&0\end{pmatrix},\qquad
 D=T(2)=\begin{pmatrix}2&1\\1&0\end{pmatrix},\qquad
 E=T(1)^2=\begin{pmatrix}2&1\\1&1\end{pmatrix}.
\]
For adjacent steps of a path, assign the block $E$ when the steps agree and
$D$ when they differ.  If the product of the $a+b-1$ adjacency blocks is
$P(\omega)$, then
\[
  M(\omega)=P(\omega)_{11}.
\]
This is exactly Schiffler's rule replacing an equal pair by the digits $1,1$
and an unequal pair by the digit $2$.

For completeness, this identity can also be obtained directly from the graph.
Expose the terminal edge of a square snake and keep two states: $p$ is the
number of perfect matchings, and $q$ is the number of matchings that leave
exactly the two terminal vertices unmatched.  Attaching one square gives

\begin{align*}
 (p,q)&\longmapsto(p+q,p)
   &&\text{for a straight continuation},\\
 (p,q)&\longmapsto(p+q,q)
   &&\text{at a turn}.
\end{align*}

Indeed, the two new vertices are matched either to each other or to the two
vertices of the shared edge.  Starting from $(2,1)$ for one square, the two
half-tile attachments associated with an adjacent pair of path steps act by
$E$ when the steps agree and by $D$ when they differ.  Thus the first state is
$P(\omega)_{11}$ and counts the actual edge subsets that meet every graph
vertex exactly once.

Close the step word cyclically.  Every path in $\cD(a,b)$ begins with $\RR$
and ends with $\UU$, so the closing edge is unequal.  Thus the cyclic matrix is
\[
  F(\omega)=D P(\omega).
\]
The number of unequal edges in the closed binary word is even.  Hence
$\det F(\omega)=1$.  For every cyclic digit cut, write the corresponding
product as
\[
  F_i=\begin{pmatrix}p_i&r_i\\q_i&s_i\end{pmatrix}.
\]
All these products have the same trace.  The two fixed points of the associated
linear fractional transformation differ by
\[
  \frac{\sqrt{(p_i-s_i)^2+4r_iq_i}}{q_i}
  =\frac{\sqrt{\tr(F)^2-4}}{q_i}.
\]
To relate this calculation to the standard Lagrange number, let $\alpha_i$ be
the positive purely periodic continued fraction at the $i$th cut.  Reversing
the period transposes its matrix product.  If $\alpha_i'$ is the negative
conjugate root, then $-1/\alpha_i'$ is the positive fixed point of the reversed
period.  Hence the backward tail is $-\alpha_i'$, and the usual forward-plus-
backward approximation coefficient is precisely
$\alpha_i-\alpha_i'$.  Taking the maximum over the finitely many cyclic cuts
therefore gives
\begin{equation}\label{eq:L-square}
  L(\omega)^2=\frac{\tr(F(\omega))^2-4}{q_{\min}(\omega)^2},
  \qquad q_{\min}(\omega)=\min_i q_i.
\end{equation}
Equation~\eqref{eq:L-square} gives an exact reduced rational comparison for
the Lagrange order, since all Lagrange values here are positive.

\section{Band-graph classes}\label{sec:band-classes}

Close the step word of $\omega$ cyclically.  Encode an equal cyclic edge by
the pair $1,1$ and an unequal edge by $2$.  Since every path begins with $\RR$
and ends with $\UU$, the closing edge contributes a leading $2$, and the
resulting cyclic coefficient word is exactly
\[
 C(\omega)=(2,A(\omega)).
\]
We regard $C(\omega)$ with its intrinsic parsing into the blocks $(1,1)$ and
$(2)$.  Equivalently, this parsing is recovered from the flattened digit word:
every cyclic run of $1$'s has even length, and every $2$ is a singleton block.

\begin{lemma}[Recovery from the cyclic code]\label{lem:code-recovery}
The cyclic equality/difference word of the steps, equivalently the cyclic
coefficient word $C(\omega)$, determines the cyclic step word of $\omega$ up
to global complementation.  If two paths in the same $\cD(a,b)$ have either
of these codes related by a rotation or reversal, then their cyclic step words
are related by a rotation or reversal, without complementation.
\end{lemma}

\begin{proof}
Pair the $1$'s in each cyclic run of $C(\omega)$ and replace every pair by an
equal edge; replace each $2$ by an unequal edge.  This recovers the cyclic
equality/difference word of the steps.  The parsing commutes with cyclic
rotation and reversal, because the blocks $(1,1)$ and $(2)$ are palindromic.
After one step is chosen, all remaining steps are forced.  The other initial
choice complements every step.  Global complementation exchanges the numbers
of $\RR$'s and $\UU$'s, and hence cannot take one word with $a$ copies of
$\RR$ and $b$ copies of $\UU$ to another because $a\ne b$.
\end{proof}

We next record the coprime cycle lemma in the form needed here.  Give an
$\RR$-step weight $b$ and a $\UU$-step weight $-a$.

\begin{lemma}[Unique admissible rotation]\label{lem:unique-rotation}
Every cyclic word with $a$ copies of $\RR$ and $b$ copies of $\UU$ has a
unique rotation in $\cD(a,b)$.
\end{lemma}

\begin{proof}
The total weight is zero.  No two proper cyclic prefix sums are equal: the
intervening proper segment would have
$b\,\#_{\RR}=a\,\#_{\UU}$, and coprimality would force it to contain all
$a$ right steps and all $b$ up steps.  Thus the cyclic prefix sums are
distinct.  Rotating immediately after their unique minimum makes every
proper prefix sum positive, while any other rotation has a negative prefix
ending at that minimum.  This is precisely the defining prefix inequality.
\end{proof}

Let $\tau(\omega)$ be the unique admissible rotation of the reversed step
word.  Lemma~\ref{lem:unique-rotation} shows at once that $\tau$ is an
involution.

\begin{lemma}[Intrinsic tile cycle]\label{lem:tile-cycle}
The abstract graph $\overline G(\omega)$ determines the cyclic
equality/difference word of the steps up to rotation and reversal.  Conversely,
a rotation or reversal of that word induces an isomorphism of the band graphs.
\end{lemma}

\begin{proof}
If $a+b<5$, then $(a,b)$ is either $(2,1)$ or $(3,1)$, and the prefix
condition forces $\cD(a,b)=\{\RR^a\UU\}$.  The assertion is therefore
immediate.  Assume henceforth that $a+b\geq5$.

Schiffler's construction places one square at every step midpoint and every
interior path vertex, adds the closing square, and glues the two terminal
edges; see \cite{CanakciSchifflerBand,Schiffler}.  Before the final gluing, the
square centers form a monotone width-one ribbon.  It contains no $2\times2$
block, so every chordless four-cycle in its one-skeleton is the boundary of a
single square.  After gluing, a new four-cycle using the glued edge would give
a three-edge path between its endpoints in the open ribbon.  The only such
paths are the other three sides of the two squares incident with that edge;
the alternative route traverses the remaining ribbon, which has at least six
tiles.  Thus the chordless four-cycles of $\overline G(\omega)$ are precisely
its $2(a+b)$ tiles, and their shared-edge intersection graph is a cycle.

The tile cycle is bipartite.  One bipartition class consists of the step
midpoint squares; on each of them the two shared edges are opposite.  The other
class consists of the path-vertex squares.  On such a square the shared edges
are opposite when the two cyclically adjacent steps agree and adjacent when
they differ.  Since the closed step word contains both letters, the latter
class contains an adjacent-edge square and is therefore distinguished from the
midpoint class by the abstract graph.  Reading it around the tile cycle
recovers the cyclic equality/difference word up to the dihedral symmetries of
that cycle.  Reversing this construction proves the converse.
\end{proof}

\begin{proof}[Proof of Theorem~\ref{thm:band-classes}]
 Lemma~\ref{lem:tile-cycle} says that the band graphs are isomorphic if and only
 if the cyclic equality/difference codes are related dihedrally.  By
 Lemma~\ref{lem:code-recovery}, their step words are then related by rotation
or reversal.  Lemma~\ref{lem:unique-rotation} leaves only $\omega$ itself in
the rotation class and $\tau(\omega)$ in the reversed class.  The converse
follows directly from the construction.
\end{proof}

The classification also gives a closed count.  Put $N=a+b$ and define
\[
 R_{a,b}=\begin{cases}
 \displaystyle\binom{(N-1)/2}{\lfloor a/2\rfloor},&N\text{ odd},\\[8pt]
 \displaystyle\binom{(N-2)/2}{(a-1)/2},&N\text{ even}.
 \end{cases}
\]
In the second case $a$ and $b$ are both odd.  The usual coprime cycle lemma
gives $|\cD(a,b)|=\binom Na/N$.  Aperiodicity follows from coprimality.  When
$N$ is odd, a reflection fixes one position and pairs the remaining positions;
when $N$ is even, only the reflections fixing two positions can preserve the
odd number $a$ of right steps.  Counting the fixed binary necklaces gives
$R_{a,b}$.  Burnside's lemma for the involution $\tau$ now yields the following.

\begin{corollary}[Number of band-graph classes]\label{cor:band-count}
The number of band-graph isomorphism classes in $\cD(a,b)$ is
\[
 \frac12\left(\frac1N\binom Na+R_{a,b}\right).
\]
For $\cD(n,n-1)$ this becomes
\[
 \frac12\left(\operatorname{Cat}_{n-1}
 +\binom{n-1}{\lfloor(n-1)/2\rfloor}\right).
\]
\end{corollary}

\section{The counterexample}\label{sec:counterexample}

Direct prefix checking places the two words of
Theorem~\ref{thm:counterexample} in $\cD(17,9)$.  Their adjacency digit words
both have $37$ digits and their cyclic words both have $38$ digits.  Exact
matrix multiplication gives, for each path,
\begin{align*}
 \tr(F)&=20034079647,\\
 \tr(F)^2-4&=401364347302339644605,\\
 q_{\min}&=6262579207.
\end{align*}
Therefore
\begin{equation}\label{eq:counter-value}
 L(\omega)^2=L(\omega')^2
 =\frac{401364347302339644605}{39219898323948748849}.
\end{equation}

It remains to separate the unlabelled band graphs.  We use an invariant that
is recovered from the abstract graph.  In each of the two band graphs, the
chordless $4$-cycles
are exactly the $52$ square tiles.  Their edge-intersection graph is a
$52$-cycle.  At each tile, record $1$ if the two edges shared with its two
neighboring tiles are adjacent, and $0$ if they are opposite.  An abstract
graph isomorphism must rotate or reverse this cyclic word.  The lexicographically
least dihedral representatives for the two graphs are
\begin{align*}
 &{\tt 0000000000010100010100010001010100010100000100010101},\\
 &{\tt 0000000000010001000101000101000001010001010100010101}.
\end{align*}
They differ, so the band graphs are not isomorphic.  This proves
Theorem~\ref{thm:counterexample}.

For reference, the corresponding least dihedral representatives of the cyclic
coefficient words are
\begin{align*}
 &{\tt 11111111112211221121122211221111211222},\\
 &{\tt 11111111112112112211221111221122211222}.
\end{align*}

\begin{proposition}[Computational minimality through length $25$]
\label{prop:minimal}
No coprime endpoint of total length at most $25$ contains a counterexample to
Problem 6.3.  Thus the example above is computationally minimal in total
length.
\end{proposition}

\begin{proof}[Exact finite verification]
For every coprime pair $a>b>0$ with $a+b\leq25$, rational-Dyck words are
generated recursively, retaining precisely the one-step extensions satisfying
the prefix inequality.  This gives all $99$ endpoints and all $934{,}635$
paths in the stated range.  For each path, equation~\eqref{eq:L-square} is
evaluated with arbitrary-precision integer matrices.  Equality of two
Lagrange squares is decided by exact cross multiplication, and the band class
is represented by the least rotation among the cyclic turn word and its
reversal.  Grouping by the exact Lagrange square produces $471{,}455$ score
classes; every class is contained in one band-graph class.

For each endpoint, the number generated is also checked against the coprime
cycle-lemma count $\binom{a+b}{a}/(a+b)$.  Hence the recursion omits no path in
the stated range.

An independent C++ implementation reproduces every endpoint, path, and
score-class count through length $25$ and again verifies that no score class
meets two band classes.  Its unsigned $128$-bit arithmetic is exact on this
range: the sum of the cyclic coefficients is at most $50$, so every continuant
is at most $\Fib_{51}$ and every squared quantity is far below $2^{128}$.  A
third exact implementation verifies the displayed length-$26$ pair and
constructs its two abstract band graphs directly.  The agreement of the two
exhaustive enumerations proves the proposition.
\end{proof}

\section{Structural matching covers}\label{sec:matching-structure}

We first prove that the parity of the matching score does not depend on the
path.  Work over $\Ftwo$ and put $F=T(1)$ and $S=T(2)$.  Direct multiplication
gives
\[
 F^3=I,\qquad E=F^2=F^{-1},\qquad SFS=F^{-1}.
\]
Write a path by its nonempty runs as
\[
 \RR^{r_1}\UU^{u_1}\cdots\RR^{r_k}\UU^{u_k}.
\]
Its adjacency product modulo two is
\[
 E^{r_1-1}SE^{u_1-1}S\cdots
 SE^{r_k-1}SE^{u_k-1}.
\]
Moving every occurrence of $S$ to the right by
$SF^j=F^{-j}S$, and using $\sum r_i=a$ and $\sum u_i=b$, reduces this product
to
\begin{equation}\label{eq:parity-normal-form}
 P(\omega)\equiv F^{b-a}S\pmod2.
\end{equation}
The upper-left entry of $F^jS$ is zero precisely when $j\equiv0\pmod3$.
This proves Theorem~\ref{thm:parity}.

\begin{corollary}[Unit-gap covers]\label{cor:unit-gap}
If $x,y\in\cD(a,b)$ and $M(y)=M(x)+2$, then $y$ covers $x$ in the matching
order.
\end{corollary}

\begin{proof}
All realized matching scores in this endpoint class have the same parity by
Theorem~\ref{thm:parity}.  No integer of that parity lies strictly between
$M(x)$ and $M(y)$.
\end{proof}

The local interchange proposed after Problem 6.5 also has a closed score
formula, although it does not have a uniform direction.  Suppose two valid
paths differ by an interior interchange
\[
 \text{old}=\alpha\RR\UU\beta,
 \qquad
 \text{new}=\alpha\UU\RR\beta.
\]
Let $\ell$ and $r$ be the steps immediately before and after the interchanged
pair.  Write the matrix products strictly before and after the three affected
adjacency blocks as
\[
 P=\begin{pmatrix}p&r_0\\q&s\end{pmatrix},
 \qquad
 Q=\begin{pmatrix}u&v\\w&z\end{pmatrix}.
\]

\begin{proposition}[Exact local-interchange formula]
\label{prop:local-swap}
The matching-score difference is given by the following four cases.
\begin{center}
\renewcommand{\arraystretch}{1.15}
\begin{tabular}{@{}c c c c@{}}
\toprule
$(\ell,r)$ & old blocks & new blocks & $M(\mathrm{old})-M(\mathrm{new})$\\
\midrule
$(\RR,\RR)$ & $EDD$ & $DDE$ & $2(r_0u-pw)$\\
$(\RR,\UU)$ & $EDE$ & $DDD$ & $2(pw+r_0u+r_0w)>0$\\
$(\UU,\RR)$ & $DDD$ & $EDE$ & $-2(pw+r_0u+r_0w)<0$\\
$(\UU,\UU)$ & $DDE$ & $EDD$ & $2(pw-r_0u)$\\
\bottomrule
\end{tabular}
\end{center}
\end{proposition}

\begin{proof}
Only the first row of $P$ and the first column of $Q$ contribute to the
upper-left entry of the complete product.  Substitution of
\[
 EDE-DDD=\begin{pmatrix}0&2\\2&2\end{pmatrix},
 \qquad
 EDD-DDE=\begin{pmatrix}0&-2\\2&0\end{pmatrix}
\]
gives the four displayed expressions.  All entries of $P$ and $Q$ that occur
in the middle two cases are nonnegative, with $p,u>0$.  The sum
$pw+r_0u+r_0w$ could vanish only if both exterior products were empty; the
entire path would then contain two right and two up steps, contrary to
$a>b$.  Hence the two displayed signs are strict.
\end{proof}

Thus the two mixed-neighbor cases force opposite directions, while the
equal-side cases compare the prefix and suffix continuant ratios.  Combined
with Corollary~\ref{cor:unit-gap}, the identities
$|r_0u-pw|=1$ or $pw+r_0u+r_0w=1$ certify a matching cover without enumerating
the remaining paths.  The formula alone is not a complete cover criterion:
larger differences can still be consecutive among the realized scores.

\section{Path-prefix bounds and all cover relations}\label{sec:covers}

Fix a feasible prefix $\alpha$ of length $k\geq1$.  Let $P_\alpha$ be the
product of its $k-1$ known adjacency blocks, and put $r=a+b-k$.  Every unknown
block $B$ satisfies $D\leq B\leq E$ entrywise.  Since all entries are
nonnegative, multiplication preserves these inequalities.

\subsection{Matching bounds}

Every completion $\omega$ of $\alpha$ satisfies
\begin{equation}\label{eq:M-bounds}
 \ell_M(\alpha):=(P_\alpha D^r)_{11}
 \leq M(\omega)\leq
 (P_\alpha E^r)_{11}=:u_M(\alpha).
\end{equation}

\subsection{Lagrange bounds}

Trace monotonicity gives
\[
 T_-(\alpha):=\tr(DP_\alpha D^r)
 \leq \tr(F(\omega))\leq
 \tr(DP_\alpha E^r)=:T_+(\alpha).
\]
At the unshifted cut,
\[
 q_{\min}(\omega)\leq q_0(\omega)
 \leq (DP_\alpha E^r)_{21}=:Q_+(\alpha).
\]
This yields the lower bound
\begin{equation}\label{eq:L-lower}
 \ell_L(\alpha)=
 \frac{\max\{0,T_-(\alpha)^2-4\}}{Q_+(\alpha)^2}
 \leq L(\omega)^2.
\end{equation}

Let $e(\alpha)$ be the number of equal adjacencies already present in the
prefix.  The completed cyclic digit word has at least $a+b+e(\alpha)$ digits.
At any cut, its lower-left entry is a positive-digit continuant of all digits
except the first.  It is therefore at least the all-ones value
$\Fib_{a+b+e(\alpha)}$, where $\Fib_0=0$, $\Fib_1=1$.  Hence
\begin{equation}\label{eq:L-upper}
 L(\omega)^2\leq
 u_L(\alpha):=
 \frac{\max\{0,T_+(\alpha)^2-4\}}
 {\Fib_{a+b+e(\alpha)}^2}.
\end{equation}
At a leaf, all four bounds are replaced by the corresponding exact score.

\subsection{Best-first level theorem}

The feasible prefixes form a finite rooted tree.  Maintain a frontier
antichain partitioning the not-yet-emitted leaves.  For the matching traversal,
key a nonleaf node by $\ell_M$; for the Lagrange traversal, key it by
$\ell_L$.  Key a leaf by its exact matching score or exact reduced Lagrange
square.  Repeatedly expand a minimum-key nonleaf node or emit a minimum-key
leaf, with nonleaf nodes taking precedence at an equal key.

\begin{lemma}\label{lem:best-first}
For either choice of score, leaves are emitted in nondecreasing exact score
order, and every equality level is emitted as one complete batch.
\end{lemma}

\begin{proof}
Suppose a leaf of score $s$ is about to be emitted while an unemitted leaf of
score $t<s$ remains.  Exactly one current frontier prefix contains that leaf.
Its key is a lower bound for all its completions, so its key is at most $t$ and
it has priority over the leaf of key $s$, a contradiction.  The same argument
with $t=s$, together with nonleaf-first tie handling, proves that the equality
batch is complete before a larger leaf is emitted.
\end{proof}

Lemma~\ref{lem:best-first} proves the first statement of
Theorem~\ref{thm:global-cover}.  Indeed, in a strict order pulled back from a
scalar score, a path in level $A$ is covered by a path in level $B$ exactly
when $A$ and $B$ are consecutive distinct scalar levels.

\subsection{A direct cover certificate}

For a prefix $\alpha$, let $I(\alpha)=[\ell(\alpha),u(\alpha)]$ be the
applicable interval from \eqref{eq:M-bounds}, or from
\eqref{eq:L-lower} and \eqref{eq:L-upper}.

\begin{proposition}[Prefix-antichain criterion]\label{prop:antichain}
Let $x,y\in\cD(a,b)$ have exact scores $s<t$ in either order.  Then $y$ covers
$x$ if and only if there is a finite prefix antichain $\mathcal A$ whose
cylinders partition $\cD(a,b)$ and such that every $\alpha\in\mathcal A$
satisfies
\[
  u(\alpha)\leq s
  \qquad\hbox{or}\qquad
  \ell(\alpha)\geq t.
\]
\end{proposition}

\begin{proof}
If the antichain exists, the score of every path lies outside $(s,t)$, so no
third level occurs between $x$ and $y$.  Conversely, refine every prefix whose
interval meets $(s,t)$.  The tree is finite.  At a leaf the interval is the
singleton exact score.  If $x<y$ is a cover, every leaf is at most $s$ or at
least $t$, so the refinement terminates with the required antichain.
\end{proof}

This proves the second statement of Theorem~\ref{thm:global-cover}.  Notice
that Proposition~\ref{prop:antichain} is a verifiable global interval-exclusion
certificate.  It does not merely recompute the two endpoint scores.

There is an equivalent deterministic formulation.  For exact scores $s<t$,
define $\mathsf C_{s,t}(\alpha)$ recursively: it is true when
$u(\alpha)\leq s$ or $\ell(\alpha)\geq t$; it is false at an otherwise
undecided leaf; and at an otherwise undecided nonleaf prefix it is the
conjunction of $\mathsf C_{s,t}$ over all feasible one-step extensions.
Then
\[
  x\lessdot y
  \quad\Longleftrightarrow\quad
  \mathsf C_{s,t}(\RR)\ \hbox{is true}.
\]
Thus the criterion does not assume a cover oracle.  A successful recursion
returns the terminal true prefixes as its antichain certificate; a failed
recursion returns an intermediate path.

\begin{remark}[Certificate size]\label{rem:scope}
The full set of leaves is the worst-case certificate in the necessity proof.
Thus Theorem~\ref{thm:global-cover} gives an exact global characterization but
does not imply a uniformly short certificate or a finite catalogue of local
swaps.  For scale,
the matching cover $X_{13}\lessdot Y_{13}$ has a certificate with $122$
terminal prefixes after $143$ expansions, while $\cD(13,12)$ has $208{,}012$
paths.
\end{remark}

\section{An unbounded-distance matching-cover family}\label{sec:nonlocal}

We give the proof of Theorem~\ref{thm:nonlocal}.  Let
\[
 P_{-1}=1,\quad P_0=0,\quad P_1=1,\qquad
 P_{j+1}=2P_j+P_{j-1}.
\]
Then
\[
 D^j=\begin{pmatrix}P_{j+1}&P_j\\P_j&P_{j-1}\end{pmatrix}.
\]
Put $m=2n-2$ and $N=m-1$.  If an adjacency-block word has initial block $E$
and exactly two further $E$ blocks at positions $1\leq i<j\leq m-1$, expansion
of $E=D+\left(\begin{smallmatrix}0&0\\0&1\end{smallmatrix}\right)$ gives
\begin{equation}\label{eq:three-E}
 M(i,j)=P_{m+1}+f(i)+f(j)+P_iP_{j-i-2}P_{m-1-j},
 \qquad f(k)=P_kP_{N-k}.
\end{equation}
The convention $P_{-1}=1$ covers adjacent positions.

Extremal continuant rearrangements such as Ramharter's
\cite{Ramharter} compare permutations of a fixed digit multiset.  Here the
paired-$1$ constraint, fixed endpoints, and rational-Dyck prefix condition
change the admissible family, so we use the following self-contained Pell
comparison.

The required Pell comparison is the following elementary lemma.

\begin{lemma}\label{lem:pell-extrema}
For odd $N\geq5$ and $1\leq k\leq N-1$,
\[
 f(k)\leq f(1),\qquad
 f(k)\geq f(2),
\]
with equality in the second inequality only at $k=2,N-2$.  Away from those
two positions, $f(k)\geq f(4)$.
\end{lemma}

\begin{proof}
Write $N=2s+1$ and use $f(k)=f(N-k)$.  For
$1\leq k\leq s-1$, the determinant identities
\[
 f(k+2)-f(k)=2(-1)^kP_{N-2k-2},\qquad
 f(k)-f(k+1)=(-1)^{k+1}P_{N-2k-1}.
\]
have positive Pell numbers on their right-hand sides.  Hence the even
subsequence increases from $f(2)$ and the odd subsequence decreases from
$f(1)$ up to the midpoint.  The second identity interlaces them: an even term
is smaller than the next odd term, while an odd term is larger than the next
even term.  This proves the stated maximum and minimum, including the equality
cases.  Every even term other than $f(2)$ is at least $f(4)$, and every
odd $k$ with $3\leq k\leq s-1$ satisfies
$f(k)>f(k+1)\geq f(4)$.  If $s$ is odd, symmetry gives
$f(s)=f(s+1)\geq f(4)$; and $f(1)\geq f(4)$ follows from the maximum
assertion.  Symmetry supplies the other half.
\end{proof}

We first verify the path conditions explicitly.  At the $j$th up-step in the
alternating part of $X_n$, where $1\leq j\leq n-3$, the numbers of right and up
steps are $j+2$ and $j$; hence the prefix margin is
\[
 (n-1)(j+2)-nj=2n-2-j>0.
\]
The two final up-steps have margins $n$ and $0$.  For $Y_n$, the $j$th up-step
in the alternating part has counts $j+1,j$ and margin $n-1-j>0$ for
$1\leq j\leq n-4$; its three remaining up-steps have margins $n+1,1,0$.
Thus both displayed words lie in $\cD(n,n-1)$.

Moreover every path in this set begins with $\RR\RR$: the alternative prefix
$\RR\UU$ would require $n\leq n-1$.  Every such path ends with $\UU$, because
deleting a final $\RR$ would leave all $n-1$ up-steps but only $n-1$ right-steps.
There are $2n-2$ adjacency blocks, while a binary word beginning in $\RR$ and
ending in $\UU$ has an odd number of unequal adjacencies.  Its block word
therefore begins with $E$ and has an odd number of $E$ blocks.

The block words of the two displayed paths are
\[
 A(X_n)=EED^{2n-5}E,\qquad
 A(Y_n)=ED^{2n-8}EDED^2.
\]
Thus
\[
 M(X_n)=P_{m+1}+f(1),\qquad
 M(Y_n)=P_{m+1}+f(2)+f(4).
\]
Their difference is positive because, with $h=m-5\geq1$,
\[
 M(Y_n)-M(X_n)=12P_h-P_{h+1}>0.
\]
Indeed $P_{h+1}=2P_h+P_{h-1}<3P_h$.

Now classify an arbitrary path by its odd number of equal adjacency blocks.
A path with one $E$ has score below
$M(X_n)$.  For a three-$E$ word ending in $E$, equation
\eqref{eq:three-E} and Lemma~\ref{lem:pell-extrema} give score at most
$M(X_n)$.  For a three-$E$ word ending in $D$, the same lemma gives score at
least $M(Y_n)$, except apparently for positions $\{2,m-3\}$.  In that case the
interaction term in \eqref{eq:three-E} gives
\[
 M(2,m-3)-M(Y_n)=2P_{m-7}>0.
\]
Finally, in a word with at least five $E$ blocks, keep the initial $E$ and any
two nonfinal $E$ blocks and replace every other $E$ by $D$.  The result is a
three-$E$ word ending in $D$, and at least one nonfinal replacement is strict.
Indeed, for nonempty products $P,Q$ on the two sides of a replaced block,
\[
 \bigl(P(E-D)Q\bigr)_{11}=P_{12}Q_{21}>0.
\]
Hence the original score is above $M(Y_n)$.  These cases
exclude the whole open interval and prove the cover assertion.  Direct
comparison of the step words gives Hamming distance $2n-6$.

It remains to establish the stronger initial-level statement.  With one
$E$-block, alternation forces the unique block word $ED^{m-1}$ and score
$P_{m+1}$.  A three-$E$ word ending in $E$ has the form
\[
 A_r=ED^{2r}ED^{m-2r-3}E,
 \qquad 0\leq r\leq n-3.
\]
Indeed, the endpoint counts force the movable $E$-position to be odd, and a
direct prefix check shows that every displayed word decodes to a path in
$\cD(n,n-1)$.  Equation~\eqref{eq:three-E} gives
\[
 M(A_r)=P_{m+1}+f(2r+1).
\]
The first determinant identity in the proof of
Lemma~\ref{lem:pell-extrema} shows that
\[
 f(1)>f(3)>\cdots>f(m-3).
\]
Thus, in decreasing order of $r$, the $A_r$ are the $n-2$ singleton levels
immediately above the one-$E$ minimum, and $A_0=X_n$ is the largest of them.

For a three-$E$ word ending in $D$, equality in the lower estimate
$M\geq M(Y_n)$ can occur only for the position pairs
$(2,4)$ and $(m-5,m-3)$.  The first decodes to a word beginning
$\RR\RR\UU\UU$ and violates the rational-Dyck inequality at that prefix.
The second decodes to $Y_n$.  The strict replacement argument already shows
that every word with at least five $E$-blocks lies above this level.  Hence
$Y_n$ is the unique next path.  This completes the proof of all assertions in
Theorem~\ref{thm:nonlocal}.

\section{Computational verification and formalization status}
\label{sec:formal}

All computations used in the paper are finite and exact.  The two paths in
Theorem~\ref{thm:counterexample} are checked by direct prefix tests, integer
matrix multiplication, and a graph-intrinsic dihedral invariant.  The bounded
claim in Proposition~\ref{prop:minimal} is checked by two independent
enumerations.  The best-first construction and representative prefix
certificates have also been compared with exhaustive enumeration on finite
ranges.  Every comparison uses integers or reduced rational squares.

The finite statement in Proposition~\ref{prop:minimal} is confined to total
length at most $25$.  In contrast, Theorems~\ref{thm:band-classes},
\ref{thm:global-cover},~\ref{thm:parity}, and~\ref{thm:nonlocal} are uniform
mathematical statements proved in the preceding sections.

The accompanying Lean~4 development uses Mathlib~\cite{Lean4,Mathlib}.  It
formalizes the rational-Dyck carrier, both exact matrix scores, and the literal
cyclic square-band quotient.  It proves that every chordless four-cycle of the
band is a tile, that the shared-edge graph of the tiles is the required cycle,
and that its adjacent-edge decoration is exactly the cyclic
equality/difference code.  This yields Theorem~\ref{thm:band-classes} on the
literal quotient graph.  Burnside counting of the resulting involution,
including the odd- and even-reflection fixed-point counts, gives
Corollary~\ref{cor:band-count} and its Catalan specialization.

For the counterexample, Lean checks the two carrier conditions, the displayed
identity for the square of the literal real periodic Lagrange value, and
nonisomorphism of the literal quotient bands.  It also
proves the minimality theorem for every coprime endpoint of total length at
most $25$.  The finite search contributes ordinary data only: the kernel
recomputes all $934{,}635$ path assignments, the $471{,}455$ exact score
classes, the $22{,}569$ prefix regions, and the complete endpoint range.

For the cover problem, the development defines the actual frontier state and
its two transitions: erase and expand a selected nonleaf by precisely its
feasible children, or erase and append a selected leaf to the output.  It
proves preservation of the frontier partition, antichain, score-order, and
completeness invariants.  The potential
$\sum_{\alpha}3^{a+b-|\alpha|}$ strictly decreases at every transition, so the
traversal terminates with the output a permutation of all carrier paths.  This
constructs complete best-first runs for both manuscript keys and proves their
score order and equality-batch completeness.  The development also proves the
consecutive-level and complete-bipartite
characterizations, both prefix-antichain equivalences, the matching and
Lagrange prefix bounds, the endpoint-parity and unit-gap results, all four
local-interchange identities, the Pell extremal lemma, and the complete
nonlocal family of Theorem~\ref{thm:nonlocal}.  The large minimality
certificate is split into independently compiled class, assignment,
compact-frontier, and endpoint modules.  The publication root, the exact
manuscript formula map, the theorem map, and the axiom report are compiled with
\texttt{--trust=0} under a 32GB aggregate memory limit.

The development has no project axiom, opaque theorem, native-evaluation
bridge, or published-result premise.  It defines a square snake as an explicit
finite vertex-and-edge type and a perfect matching as an edge selection meeting
every vertex exactly once; local equivalences prove the two-state recurrence
above.  It also defines the periodic Lagrange value by the standard maximum of
forward-plus-backward cut values, proves that the backward value is the
negative conjugate, and derives equation~\eqref{eq:L-square}.  The axiom report
for every publication endpoint contains only Lean/Mathlib's standard
\texttt{propext}, \texttt{Classical.choice}, and \texttt{Quot.sound}.

No theorem-valued certificate or search conclusion is assumed.  The same is
true of every path-code translation, graph isomorphism, cover implication,
counting identity, and final composition.  Every numbered result has an
explicit entry in the accompanying manuscript-to-Lean map.  The
reproducibility documentation records the exact commands, generated-source
counts, formula and theorem maps, and trust report.

\paragraph{Data and code availability.}
The Lean source, exact certificates, verification scripts, and current
manuscript PDF are available at
\url{https://github.com/crabsatellite/lattice-path-orders}; versioned releases
are archived through Zenodo.

\section*{Declaration of generative AI and AI-assisted technologies in the
manuscript preparation process}

During the preparation of this work, the author used an author-developed local
implementation of the Proof Engine framework for AI-assisted mathematical
research.  The implementation operates through replaceable AI-agent and
harness components; the present work used OpenAI Codex within that
implementation.  The framework is described in \emph{Proof Engine
Infrastructure}~\cite{LiProofEngine2026}.  The author reviewed and edited the
resulting material and assumes full responsibility for the definitions,
statements, proofs, citations, code, and exposition.

\section*{Statements and declarations}

\paragraph{Competing interests.}
The author declares no known competing financial or non-financial interests
that are directly or indirectly related to this work.

\paragraph{Funding.}
No funding was received for conducting this study or preparing the manuscript.

\section*{Acknowledgements}

The author thanks the developers of Lean and Mathlib for the formalization
environment.

\bibliographystyle{amsplain}
\bibliography{references}

\end{document}
